\begin{CNbox}[unbreakable]{obj}{Learning Objectives}

By the end of this module, you will be able to:

\begin{itemize}
\tightlist
\item
  Explain the complete GSM architecture and the role of each network element
\item
  Describe how GSM handles mobility, location management, and authentication
\item
  Articulate \textbf{why} GPRS was introduced and what limitations of GSM it addressed
\item
  Explain the GPRS packet-switched overlay architecture
\item
  Describe PDP context activation and GPRS mobility management
\item
  Explain how EDGE improved data rates through modulation changes
\item
  Trace the evolution path from GSM through EDGE to 3G
\end{itemize}

\end{CNbox}

\Needspace{14\baselineskip}

\CNsection{1}{GSM Architecture}

\CNchained\CNsubsection{1.1}{WHY GSM?}

Before GSM (Global System for Mobile Communications), Europe had multiple incompatible analog cellular systems (NMT, TACS, C-Netz). A subscriber in France could not roam to Germany. GSM was created to solve three fundamental problems:

\begin{enumerate}
\def\labelenumi{\arabic{enumi}.}
\tightlist
\item
  \textbf{Interoperability} — One standard across all of Europe (and eventually the world)
\item
  \textbf{Security} — Analog systems had no encryption; anyone with a scanner could listen
\item
  \textbf{Capacity} — Digital modulation serves more users per MHz of spectrum
\end{enumerate}

\CNkeepfig{m01-gsm-architecture}{3}

\CNsubsection{1.2}{GSM Network Architecture Diagram}

\CNfigure{m01-gsm-architecture}{GSM network architecture: MS, BSS, NSS and external networks with their interfaces}

\Needspace{15\baselineskip}

\CNsubsection{1.3}{Understanding the Architecture Layers}

The GSM architecture is organized into three logical subsystems, each solving a distinct problem:

{\def\LTcaptype{none} % do not increment counter
\Needspace{9\baselineskip}
\begin{longtable}[c]{@{}
  >{\raggedright\arraybackslash\hspace{0pt}}m{(\linewidth - 4\tabcolsep) * \real{0.2750}}
  >{\raggedright\arraybackslash\hspace{0pt}}m{(\linewidth - 4\tabcolsep) * \real{0.4750}}
  >{\raggedright\arraybackslash\hspace{0pt}}m{(\linewidth - 4\tabcolsep) * \real{0.2500}}@{}}
\toprule\noalign{}
\rowcolor{CNink}\CNth{Subsystem} & \CNth{Problem It Solves} & \CNth{Elements} \\
\CNheadrulesep
\endhead
\bottomrule\noalign{}
\endlastfoot
\textbf{BSS} (Base Station Subsystem) & Radio access — how does the phone talk to the network? & BTS, BSC \\
\textbf{NSS} (Network Switching Subsystem) & Call routing — how do calls get connected? & MSC, GMSC, HLR, VLR, AuC, EIR \\
\textbf{OSS} (Operation Support Subsystem) & Management — how do operators monitor and configure the network? & OMC, NMC \\
\end{longtable}
}

\Needspace{14\baselineskip}

\CNsection{2}{GSM Network Elements — Detailed Explanation}

\CNchained\CNsubsection{2.1}{Mobile Station (MS)}

\textbf{What it is:} The subscriber’s device, consisting of two parts:

\begin{itemize}
\tightlist
\item
  \textbf{ME (Mobile Equipment):} The physical handset hardware, identified by its IMEI
\item
  \textbf{SIM (Subscriber Identity Module):} A smart card carrying the subscriber’s identity (IMSI), authentication key (Ki), and algorithms
\end{itemize}

\textbf{Why it’s split into ME + SIM:} This separation means you can swap your SIM into any phone and keep your number, or an operator can block a stolen phone (by IMEI) without affecting the subscriber’s account.

\CNsubsection{2.2}{Base Transceiver Station (BTS)}

\textbf{What it is:} The radio equipment (antennas, transceivers, amplifiers) at a cell site.

\textbf{Why it exists:} Someone has to convert between the wired network and radio waves. The BTS handles:

\begin{itemize}
\tightlist
\item
  Transmission and reception of radio signals
\item
  Channel coding and encryption at the air interface
\item
  Frequency hopping (if configured)
\item
  Timing advance measurements (to compensate for propagation delay)
\end{itemize}

\begin{CNinsight}

\textbf{Key insight:} The BTS is intentionally ``dumb'' — it has minimal intelligence. This is a deliberate design choice so that expensive radio sites don’t need software upgrades for new features.

\end{CNinsight}

\CNsubsection{2.3}{Base Station Controller (BSC)}

\textbf{What it is:} The intelligence behind multiple BTSs. One BSC typically controls dozens to hundreds of BTSs.

\textbf{Why it exists:} The BSC solves the \textbf{resource management} problem:

\begin{itemize}
\tightlist
\item
  \textbf{Handover decisions:} When should a call be moved from one BTS to another? The BSC collects measurements from the MS and neighboring BTSs to decide.
\item
  \textbf{Radio resource allocation:} Which frequency and timeslot should be assigned to a new call?
\item
  \textbf{Power control:} Instructs the MS and BTS to adjust power (saves battery, reduces interference)
\item
  \textbf{Transcoding:} Converts between the compressed speech codec on the air interface (13~kbps) and standard 64~kbps PCM used in the core network
\end{itemize}

\CNsubsection{2.4}{Mobile Switching Center (MSC)}

\textbf{What it is:} A telephone exchange specifically designed for mobile subscribers.

\textbf{Why it exists:} The MSC solves the \textbf{call routing for mobile users} problem. Unlike a fixed-line exchange, the MSC must handle:

\begin{itemize}
\tightlist
\item
  \textbf{Call setup and teardown} for users who may be moving
\item
  \textbf{Handover between BSCs} (inter-BSC handover)
\item
  \textbf{Interface to external networks} (PSTN, ISDN, other mobile networks)
\item
  \textbf{Supplementary services} (call forwarding, call waiting, conference calls)
\item
  \textbf{Coordination with VLR} for subscriber validation
\end{itemize}

\CNsubsection{2.5}{Gateway MSC (GMSC)}

\textbf{What it is:} A specialized MSC that acts as the entry point for calls coming FROM external networks.

\textbf{Why it exists:} When someone on a landline calls a mobile subscriber, the call arrives at the GMSC. The GMSC’s job is to \textbf{find the subscriber} — it queries the HLR to determine which MSC/\allowbreak{}VLR area the subscriber is currently in, then routes the call there. Without the GMSC, external networks would need to know the internal structure of the mobile network.

\CNsubsection{2.6}{Home Location Register (HLR)}

\textbf{What it is:} The central database containing ALL permanent subscriber information.

\textbf{Why it exists:} The HLR solves the \textbf{``where is this subscriber?''} problem. It stores:

\begin{itemize}
\tightlist
\item
  Subscriber identity (IMSI, MSISDN)
\item
  Subscription information (which services are allowed)
\item
  \textbf{Current location} (which VLR/\allowbreak{}MSC area the subscriber is registered in)
\item
  Call forwarding settings
\item
  Supplementary service data
\end{itemize}

\begin{CNinsight}

\textbf{Key insight:} There is logically ONE HLR per network (though physically it may be distributed). It is the single source of truth about every subscriber.

\end{CNinsight}

\CNsubsection{2.7}{Visitor Location Register (VLR)}

\textbf{What it is:} A temporary database associated with each MSC, holding information about subscribers currently in that MSC’s area.

\textbf{Why it exists:} The VLR solves the \textbf{performance} problem. Without it, every single call setup would require a query to the HLR (which could be far away). Instead:

\begin{enumerate}
\def\labelenumi{\arabic{enumi}.}
\tightlist
\item
  When a subscriber enters an MSC area, the VLR fetches their profile from the HLR
\item
  Subsequent operations use the local VLR copy
\item
  This reduces signaling load on the HLR and reduces call setup latency
\end{enumerate}

The VLR also stores temporary data: the TMSI (temporary identity), location area, and authentication triplets.

\CNsubsection{2.8}{Authentication Center (AuC)}

\textbf{What it is:} A protected database storing the secret authentication keys (Ki) for every subscriber.

\textbf{Why it exists:} The AuC solves the \textbf{``how do we prove this subscriber is genuine?''} problem. It:

\begin{itemize}
\tightlist
\item
  Stores each subscriber’s Ki (which never leaves the AuC or SIM)
\item
  Generates authentication triplets (RAND, SRES, Kc) on request
\item
  Protects against cloning and unauthorized access
\end{itemize}

\textbf{Why it’s separate from the HLR:} Security isolation. The Ki values are extremely sensitive. By placing them in a physically secured, separate element, a compromise of the HLR doesn’t expose authentication keys.

\CNsubsection{2.9}{Equipment Identity Register (EIR)}

\textbf{What it is:} A database of mobile equipment identities (IMEIs) classified into three lists:

\begin{itemize}
\tightlist
\item
  \textbf{White list:} Approved equipment
\item
  \textbf{Grey list:} Equipment under observation (perhaps reported lost but not confirmed)
\item
  \textbf{Black list:} Stolen or non-type-approved equipment — service denied
\end{itemize}

\textbf{Why it exists:} The EIR solves the \textbf{stolen phone} problem. Even if a thief puts a new SIM in a stolen phone, the network can identify and block the hardware by its IMEI.

\Needspace{14\baselineskip}

\CNsection{3}{GSM Interfaces}

\CNchained\CNsubsection{}{WHY standardized interfaces matter}

GSM’s genius was defining \textbf{open interfaces} between elements. This means a BTS from Ericsson can connect to a BSC from Nokia. Without standardized interfaces, operators would be locked to a single vendor for everything.

\Needspace{25\baselineskip}

\CNsubsection{3.1}{Interface Summary}

{\def\LTcaptype{none} % do not increment counter
\Needspace{21\baselineskip}
\begin{longtable}[c]{@{}
  >{\raggedright\arraybackslash\hspace{0pt}}m{(\linewidth - 6\tabcolsep) * \real{0.2500}}
  >{\raggedright\arraybackslash\hspace{0pt}}m{(\linewidth - 6\tabcolsep) * \real{0.2045}}
  >{\raggedright\arraybackslash\hspace{0pt}}m{(\linewidth - 6\tabcolsep) * \real{0.3409}}
  >{\raggedright\arraybackslash\hspace{0pt}}m{(\linewidth - 6\tabcolsep) * \real{0.2045}}@{}}
\toprule\noalign{}
\rowcolor{CNink}\CNth{Interface} & \CNth{Between} & \CNth{Protocol Stack} & \CNth{Purpose} \\
\CNheadrulesep
\endhead
\bottomrule\noalign{}
\endlastfoot
\textbf{Um} & MS \ensuremath{\leftrightarrow} BTS & LAPDm /\allowbreak{} RR /\allowbreak{} MM /\allowbreak{} CC & Air interface — the radio link \\
\textbf{Abis} & BTS \ensuremath{\leftrightarrow} BSC & LAPD over PCM (E1/\allowbreak{}T1) & Carries traffic + signaling from radio sites to controller \\
\textbf{A} & BSC \ensuremath{\leftrightarrow} MSC & MTP /\allowbreak{} SCCP /\allowbreak{} BSSAP & Connects radio subsystem to switching subsystem \\
\textbf{B} & MSC \ensuremath{\leftrightarrow} VLR & MAP & MSC queries local VLR for subscriber data \\
\textbf{C} & GMSC \ensuremath{\leftrightarrow} HLR & MAP & GMSC asks HLR where to route incoming calls \\
\textbf{D} & VLR \ensuremath{\leftrightarrow} HLR & MAP & Location updates, subscriber data transfer \\
\textbf{E} & MSC \ensuremath{\leftrightarrow} MSC & MAP /\allowbreak{} ISUP & Inter-MSC handover \\
\textbf{F} & MSC \ensuremath{\leftrightarrow} EIR & MAP & IMEI checking \\
\textbf{G} & VLR \ensuremath{\leftrightarrow} VLR & MAP & Subscriber info transfer during inter-VLR moves \\
\textbf{H} & HLR \ensuremath{\leftrightarrow} AuC & Proprietary & Authentication vector requests \\
\end{longtable}
}

\CNsubsection{3.2}{The Um Interface (Air Interface) — In Depth}

The Um interface is the most complex because it must deal with the hostile radio environment:

\begin{itemize}
\tightlist
\item
  \textbf{Multiple access:} TDMA with 8 timeslots per 200~kHz carrier (FDMA/\allowbreak{}TDMA hybrid)
\item
  \textbf{Modulation:} GMSK (Gaussian Minimum Shift Keying) — chosen for constant envelope (efficient power amplifiers) and spectral efficiency
\item
  \textbf{Frame structure:} Hyperframe \ensuremath{\to} Superframe \ensuremath{\to} Multiframe \ensuremath{\to} Frame \ensuremath{\to} Timeslot \ensuremath{\to} Burst
\item
  \textbf{Logical channels:}

  \begin{itemize}
  \tightlist
  \item
    \textbf{TCH (Traffic Channel):} Carries user speech/\allowbreak{}data
  \item
    \textbf{BCCH (Broadcast Control Channel):} System info broadcast to all phones
  \item
    \textbf{RACH (Random Access Channel):} MS requests access (slotted ALOHA)
  \item
    \textbf{PCH (Paging Channel):} Network pages MS for incoming call
  \item
    \textbf{SDCCH (Standalone Dedicated Control Channel):} Signaling (SMS, location update)
  \item
    \textbf{SACCH/\allowbreak{}FACCH:} Associated control for measurements and handover commands
  \end{itemize}
\end{itemize}

\CNsubsection{3.3}{The Abis Interface}

\textbf{Why it exists:} BTS sites are remote (mountain tops, rooftops). The Abis interface lets them connect back to the BSC over standard 2~Mbps E1 lines (or microwave links). It carries:

\begin{itemize}
\tightlist
\item
  16~kbps traffic channels (compressed speech)
\item
  Signaling messages (LAPD protocol, similar to ISDN)
\item
  Operation \& Maintenance messages
\end{itemize}

\CNsubsection{3.4}{The A Interface}

\textbf{Why it exists:} This is where the radio world meets the telephony world. The A interface carries:

\begin{itemize}
\tightlist
\item
  64~kbps PCM speech (after transcoding from 13~kbps)
\item
  BSSAP signaling (BSS Application Part) over SS7
\end{itemize}

\Needspace{14\baselineskip}

\CNsection{4}{GSM Mobility \& Location Management}

\CNchained\CNsubsection{4.1}{WHY Location Management Is Needed}

The fundamental challenge of mobile networks: \textbf{how do you deliver a call to someone whose location is unknown?}

Naive solution: Page every cell in the network. Problem: This doesn’t scale — millions of subscribers would require constant paging on every cell.

GSM’s solution: \textbf{Location Areas (LA)}

\CNsubsection{4.2}{Location Area Concept}

\begin{itemize}
\tightlist
\item
  The network is divided into \textbf{Location Areas}, each containing multiple cells (BTSs)
\item
  Each LA has a unique \textbf{LAI} (Location Area Identity) = MCC + MNC + LAC
\item
  The network only knows a subscriber’s location to the \textbf{LA granularity}
\item
  When a call arrives, the subscriber is paged in ALL cells of their current LA
\end{itemize}

\textbf{Trade-off:}

\begin{itemize}
\tightlist
\item
  Smaller LAs \ensuremath{\to} less paging load, but more location updates (signaling)
\item
  Larger LAs \ensuremath{\to} fewer location updates, but more paging traffic
\item
  Operators tune LA sizes based on traffic patterns
\end{itemize}

\CNsubsection{4.3}{Location Update Procedure}

A mobile performs a \textbf{location update} when:

\begin{enumerate}
\def\labelenumi{\arabic{enumi}.}
\tightlist
\item
  \textbf{Normal location update:} MS detects it has entered a new LA (different LAI on BCCH)
\item
  \textbf{Periodic location update:} Timer expires — proves the MS is still reachable
\item
  \textbf{IMSI attach:} MS is powered on and registers with the network
\end{enumerate}

\textbf{Location Update Flow:}

\begin{enumerate}
\def\labelenumi{\arabic{enumi}.}
\tightlist
\item
  MS sends LOCATION UPDATING REQUEST to new MSC/\allowbreak{}VLR
\item
  New VLR contacts HLR: ``Subscriber X is now here''
\item
  HLR sends subscriber profile to new VLR
\item
  HLR tells OLD VLR to delete the subscriber record
\item
  New VLR confirms to MS: LOCATION UPDATING ACCEPT
\item
  MS stores new TMSI (if reassigned)
\end{enumerate}

\CNsubsection{4.4}{TMSI — Temporary Mobile Subscriber Identity}

\textbf{Why:} Broadcasting the IMSI over the air interface would let eavesdroppers track subscribers. The TMSI:

\begin{itemize}
\tightlist
\item
  Is assigned by the VLR and has only local significance
\item
  Is changed frequently (at each location update or call)
\item
  Means the permanent identity (IMSI) rarely travels over the air
\end{itemize}

\Needspace{14\baselineskip}

\CNsection{5}{GSM Authentication \& Security}

\CNchained\CNsubsection{5.1}{WHY Authentication Is Critical}

Without authentication:

\begin{itemize}
\tightlist
\item
  Cloned SIMs could make calls billed to someone else
\item
  Unauthorized users could access the network
\item
  Eavesdroppers could intercept calls
\end{itemize}

\CNsubsection{5.2}{The Authentication Triplet Mechanism}

GSM uses a \textbf{challenge-response} protocol based on shared secret cryptography.

\textbf{The players:}

\begin{itemize}
\tightlist
\item
  \textbf{Ki:} A 128-bit secret key, stored ONLY on the SIM and in the AuC. It \textbf{never} travels over any interface.
\item
  \textbf{RAND:} A 128-bit random number generated by the AuC (the challenge)
\item
  \textbf{SRES:} A 32-bit Signed Response (the expected answer)
\item
  \textbf{Kc:} A 64-bit cipher key for encrypting the air interface
\end{itemize}

\CNkeepfig{m01-a3-a8}{2}

\textbf{How it works:}

\CNfigure{m01-a3-a8}{A3 and A8 run on both sides: the AuC and the SIM derive the same SRES and Kc from Ki and RAND}

\textbf{Step-by-step:}

\begin{enumerate}
\def\labelenumi{\arabic{enumi}.}
\tightlist
\item
  AuC generates RAND, computes SRES = A3(Ki, RAND) and Kc = A8(Ki, RAND)
\item
  The triplet (RAND, SRES, Kc) is sent to the VLR (pre-computed, stored for later use)
\item
  When authentication is needed, VLR sends RAND to the MS
\item
  SIM computes SRES = A3(Ki, RAND) using its stored Ki
\item
  MS sends SRES back to the network
\item
  VLR compares received SRES with stored SRES — if they match, subscriber is authentic
\item
  Both sides now have Kc — used with algorithm A5 to encrypt the air interface
\end{enumerate}

\textbf{Why triplets are pre-computed:} The AuC may be geographically distant. Pre-computing multiple triplets and caching them at the VLR means authentication can happen locally without real-time AuC access.

\CNsubsection{5.3}{Encryption (A5 Algorithm)}

After authentication, both the MS and BTS share Kc. The A5 algorithm generates a keystream that XORs with the traffic. This provides:

\begin{itemize}
\tightlist
\item
  \textbf{Confidentiality} on the air interface (Um only — not end-to-end)
\item
  Protection against casual eavesdropping
\end{itemize}

\textbf{Known weakness:} A5/\allowbreak{}1 has been broken. Modern networks use A5/\allowbreak{}3 (KASUMI-based) or have moved to 3G/\allowbreak{}4G with stronger security.

\Needspace{14\baselineskip}

\Needspace{19\baselineskip}

\CNsection{6}{WHY GPRS Was Needed}

\CNchained\CNsubsection{6.1}{The GSM Data Problem}

GSM was designed for \textbf{voice}. Data was an afterthought:

{\def\LTcaptype{none} % do not increment counter
\Needspace{8\baselineskip}
\begin{longtable}[c]{@{}
  >{\raggedright\arraybackslash\hspace{0pt}}m{(\linewidth - 4\tabcolsep) * \real{0.5152}}
  >{\raggedright\arraybackslash\hspace{0pt}}m{(\linewidth - 4\tabcolsep) * \real{0.2121}}
  >{\raggedright\arraybackslash\hspace{0pt}}m{(\linewidth - 4\tabcolsep) * \real{0.2727}}@{}}
\toprule\noalign{}
\rowcolor{CNink}\CNth{GSM Data Service} & \CNth{Speed} & \CNth{Problem} \\
\CNheadrulesep
\endhead
\bottomrule\noalign{}
\endlastfoot
CSD (Circuit-Switched Data) & 9.6~kbps & Occupies a full timeslot for entire session \\
HSCSD (High-Speed CSD) & Up to 57.6~kbps & Occupies MULTIPLE timeslots — extremely expensive \\
\end{longtable}
}

\textbf{The fundamental issue: Circuit switching is wrong for data.}

Why? Consider web browsing:

\begin{itemize}
\tightlist
\item
  You request a page (1 second of activity)
\item
  You read it (30 seconds of inactivity)
\item
  You click a link (1 second of activity)
\item
  You read again (30 seconds of inactivity)
\end{itemize}

With circuit switching, you’re \textbf{paying for and consuming a timeslot during those 30 seconds of reading}. The channel is allocated but idle 95\% of the time.

\CNsubsection{6.2}{What GPRS Provides}

\textbf{GPRS (General Packet Radio Service)} introduces \textbf{packet switching} to GSM:

\begin{itemize}
\tightlist
\item
  Resources are allocated \textbf{only when data is actually being sent}
\item
  Multiple users \textbf{share} timeslots statistically
\item
  Billing can be \textbf{per-megabyte} instead of per-minute
\item
  ``Always-on'' connectivity — no dial-up delay
\item
  Theoretical speeds: 8 timeslots \ensuremath{\times} 21.4~kbps = \textbf{171.2~kbps} (theoretical max)
\item
  Practical speeds: 30–80~kbps
\end{itemize}

\CNsubsection{6.3}{GPRS Design Philosophy}

GPRS was designed as an \textbf{overlay} on existing GSM infrastructure. Operators didn’t need to replace their entire network — they added new packet-switching nodes alongside the existing circuit-switched core. This was critical for business adoption.

\Needspace{14\baselineskip}

\CNsection{7}{GPRS Architecture}

\CNchained\CNsubsection{7.1}{New Network Elements}

GPRS introduces three key new elements:

\CNsubsubsection{SGSN — Serving GPRS Support Node}

\textbf{What it is:} The packet-switched equivalent of the MSC/\allowbreak{}VLR combination.

\textbf{Why it exists:} Someone needs to handle packet-switched sessions the way the MSC handles voice calls. The SGSN:

\begin{itemize}
\tightlist
\item
  Tracks the location of GPRS-attached mobiles (at Routing Area granularity)
\item
  Performs authentication and ciphering for packet sessions
\item
  Handles mobility management for packet data
\item
  Routes packets between the MS and the GGSN
\item
  Collects charging information
\end{itemize}

\CNsubsubsection{GGSN — Gateway GPRS Support Node}

\textbf{What it is:} The packet-switched equivalent of the GMSC — the gateway to external data networks.

\textbf{Why it exists:} The mobile network must connect to the Internet (and corporate intranets). The GGSN:

\begin{itemize}
\tightlist
\item
  Acts as the \textbf{IP router} between the GPRS network and external packet data networks (PDNs)
\item
  Assigns IP addresses to mobile devices (or relays from DHCP/\allowbreak{}RADIUS)
\item
  Performs \textbf{encapsulation/\allowbreak{}de-encapsulation} of user data using GTP (GPRS Tunneling Protocol)
\item
  Enforces policy and charging rules
\item
  From the external network’s perspective, the GGSN looks like a normal IP router
\end{itemize}

\CNsubsubsection{PCU — Packet Control Unit}

\textbf{What it is:} An addition to the BSC that handles packet-switched radio resource management.

\textbf{Why it exists:} The original BSC only understood circuit-switched resource allocation. The PCU adds:

\begin{itemize}
\tightlist
\item
  Packet scheduling on shared radio channels (PDCH — Packet Data Channels)
\item
  Segmentation and reassembly of LLC frames into radio blocks
\item
  Dynamic allocation of timeslots between voice (TCH) and data (PDCH)
\item
  Uplink state flag management and acknowledgment
\end{itemize}

\CNkeepfig{m01-gprs-data-path}{3}

\CNsubsection{7.2}{GPRS Packet Data Path}

\CNfigure{m01-gprs-data-path}{GPRS packet data path from the MS to the external packet network}

\Needspace{24\baselineskip}

\CNsubsection{7.3}{GPRS Protocol Stack}

The GPRS protocol stack introduces new layers:

{\def\LTcaptype{none} % do not increment counter
\Needspace{18\baselineskip}
\begin{longtable}[c]{@{}
  >{\raggedright\arraybackslash\hspace{0pt}}m{(\linewidth - 4\tabcolsep) * \real{0.2692}}
  >{\raggedright\arraybackslash\hspace{0pt}}m{(\linewidth - 4\tabcolsep) * \real{0.3846}}
  >{\raggedright\arraybackslash\hspace{0pt}}m{(\linewidth - 4\tabcolsep) * \real{0.3462}}@{}}
\toprule\noalign{}
\rowcolor{CNink}\CNth{Layer} & \CNth{Protocol} & \CNth{Purpose} \\
\CNheadrulesep
\endhead
\bottomrule\noalign{}
\endlastfoot
Application & HTTP, SMTP, etc. & User application \\
Transport & TCP /\allowbreak{} UDP & End-to-end reliability \\
Network & IP & Addressing and routing \\
Tunneling & GTP (GPRS Tunneling Protocol) & Encapsulates IP packets between SGSN and GGSN \\
Link & LLC (Logical Link Control) & Reliable link between MS and SGSN \\
Link & SNDCP (SubNetwork Dependent Convergence Protocol) & Header compression, segmentation \\
Radio & RLC/\allowbreak{}MAC & Radio block transmission, ARQ \\
Physical & GMSK modulation on PDCH & Bits on the air \\
\end{longtable}
}

\Needspace{18\baselineskip}

\CNsubsection{7.4}{GPRS Interfaces}

{\def\LTcaptype{none} % do not increment counter
\Needspace{14\baselineskip}
\begin{longtable}[c]{@{}
  >{\raggedright\arraybackslash\hspace{0pt}}m{(\linewidth - 6\tabcolsep) * \real{0.2750}}
  >{\raggedright\arraybackslash\hspace{0pt}}m{(\linewidth - 6\tabcolsep) * \real{0.2250}}
  >{\raggedright\arraybackslash\hspace{0pt}}m{(\linewidth - 6\tabcolsep) * \real{0.2750}}
  >{\raggedright\arraybackslash\hspace{0pt}}m{(\linewidth - 6\tabcolsep) * \real{0.2250}}@{}}
\toprule\noalign{}
\rowcolor{CNink}\CNth{Interface} & \CNth{Between} & \CNth{Transport} & \CNth{Purpose} \\
\CNheadrulesep
\endhead
\bottomrule\noalign{}
\endlastfoot
\textbf{Gb} & PCU/\allowbreak{}BSS \ensuremath{\leftrightarrow} SGSN & Frame Relay or IP & Packet data from radio to core \\
\textbf{Gn} & SGSN \ensuremath{\leftrightarrow} GGSN (same PLMN) & GTP/\allowbreak{}UDP/\allowbreak{}IP & User data tunneling within network \\
\textbf{Gp} & SGSN \ensuremath{\leftrightarrow} GGSN (different PLMN) & GTP + security & Roaming data \\
\textbf{Gi} & GGSN \ensuremath{\leftrightarrow} External PDN & IP & Connection to Internet \\
\textbf{Gr} & SGSN \ensuremath{\leftrightarrow} HLR & MAP & Subscriber data for GPRS \\
\textbf{Gs} & SGSN \ensuremath{\leftrightarrow} MSC/\allowbreak{}VLR & BSSAP+ & Coordination of CS and PS \\
\end{longtable}
}

\Needspace{14\baselineskip}

\CNsection{8}{PDP Context}

\CNchained\CNsubsection{8.1}{WHAT Is a PDP Context?}

A \textbf{PDP Context (Packet Data Protocol Context)} is the ``session'' that must exist before a mobile can send or receive packet data. Think of it as the GPRS equivalent of ``dialing in'' — but without wasting resources when idle.

A PDP context defines:

\begin{itemize}
\tightlist
\item
  \textbf{PDP Type:} Usually IPv4 or IPv6
\item
  \textbf{PDP Address:} The IP address assigned to the MS
\item
  \textbf{Access Point Name (APN):} Which external network to connect to (e.g., ``internet'', ``corporate.vpn'')
\item
  \textbf{QoS Profile:} Requested quality of service (priority, throughput, delay class)
\item
  \textbf{GGSN Address:} Which GGSN serves this context
\end{itemize}

\CNsubsection{8.2}{WHY PDP Context Exists}

Without some form of session establishment:

\begin{itemize}
\tightlist
\item
  The GGSN wouldn’t know where to send downlink packets for a given IP address
\item
  The network couldn’t enforce QoS or charging policies
\item
  There would be no tunnel established between SGSN and GGSN
\end{itemize}

\CNkeepfig{m01-pdp-activation}{3}

\CNsubsection{8.3}{PDP Context Activation Procedure}

\CNfigure{m01-pdp-activation}{PDP context activation and the first IP packets through the GTP tunnel}

\Needspace{13\baselineskip}

\CNsubsection{8.4}{PDP Context States}

{\def\LTcaptype{none} % do not increment counter
\Needspace{9\baselineskip}
\begin{longtable}[c]{@{}
  >{\raggedright\arraybackslash\hspace{0pt}}m{(\linewidth - 4\tabcolsep) * \real{0.2258}}
  >{\raggedright\arraybackslash\hspace{0pt}}m{(\linewidth - 4\tabcolsep) * \real{0.2903}}
  >{\raggedright\arraybackslash\hspace{0pt}}m{(\linewidth - 4\tabcolsep) * \real{0.4839}}@{}}
\toprule\noalign{}
\rowcolor{CNink}\CNth{State} & \CNth{Meaning} & \CNth{Resources Used} \\
\CNheadrulesep
\endhead
\bottomrule\noalign{}
\endlastfoot
\textbf{INACTIVE} & No PDP context — cannot send data & None \\
\textbf{ACTIVE} & Full context exists, data can flow & SGSN/\allowbreak{}GGSN tunnel, IP address allocated \\
\textbf{STANDBY} (optional) & Context exists but radio resources released & Tunnel maintained, no radio allocation \\
\end{longtable}
}

\Needspace{14\baselineskip}

\CNsection{9}{GPRS Mobility Management}

\CNchained\CNsubsection{9.1}{WHY GPRS Needs Its Own Mobility Management}

GSM already has location management (Location Areas, location updates). Why does GPRS need a separate system?

\textbf{Because packet-switched and circuit-switched traffic have different needs:}

\begin{itemize}
\tightlist
\item
  For a \textbf{voice call}, the network needs to find you immediately (paging must be fast)
\item
  For \textbf{packet data}, there’s more tolerance for delay — but the network needs to know which SGSN to route your downlink packets to
\end{itemize}

GPRS introduces the \textbf{Routing Area (RA)} — a subdivision of the Location Area.

\CNkeepfig{m01-la-ra}{3}

\CNsubsection{9.2}{Routing Area vs. Location Area}

\CNfigure{m01-la-ra}{A Location Area is subdivided into Routing Areas}

\begin{itemize}
\tightlist
\item
  A \textbf{Routing Area} is always a subset of (or equal to) a Location Area
\item
  RA is identified by \textbf{RAI} = LAI + RAC (Routing Area Code)
\item
  Finer granularity than LA \ensuremath{\to} more precise location for packet routing
\end{itemize}

\Needspace{13\baselineskip}

\CNsubsection{9.3}{GPRS Mobility Management States}

{\def\LTcaptype{none} % do not increment counter
\Needspace{9\baselineskip}
\begin{longtable}[c]{@{}
  >{\raggedright\arraybackslash\hspace{0pt}}m{(\linewidth - 4\tabcolsep) * \real{0.2000}}
  >{\raggedright\arraybackslash\hspace{0pt}}m{(\linewidth - 4\tabcolsep) * \real{0.6286}}
  >{\raggedright\arraybackslash\hspace{0pt}}m{(\linewidth - 4\tabcolsep) * \real{0.1714}}@{}}
\toprule\noalign{}
\rowcolor{CNink}\CNth{State} & \CNth{What the network knows} & \CNth{When} \\
\CNheadrulesep
\endhead
\bottomrule\noalign{}
\endlastfoot
\textbf{IDLE} & Nothing — MS is not GPRS-attached & Power off or GPRS detached \\
\textbf{STANDBY} & Current Routing Area & Attached but no active data transfer \\
\textbf{READY} & Current Cell (exact) & Actively sending/\allowbreak{}receiving data \\
\end{longtable}
}

\textbf{State transitions:}

\begin{itemize}
\tightlist
\item
  IDLE \ensuremath{\to} STANDBY: GPRS Attach
\item
  STANDBY \ensuremath{\to} READY: Data transfer begins (or signaling)
\item
  READY \ensuremath{\to} STANDBY: READY timer expires (no data for \textasciitilde{}44 seconds)
\item
  STANDBY \ensuremath{\to} IDLE: GPRS Detach
\end{itemize}

\CNsubsection{9.4}{Routing Area Update (RAU)}

Similar to GSM’s location update, but for the packet domain:

\textbf{When triggered:}

\begin{enumerate}
\def\labelenumi{\arabic{enumi}.}
\tightlist
\item
  MS moves to a new Routing Area (different RAI on BCCH)
\item
  Periodic RAU timer expires
\item
  Combined RA/\allowbreak{}LA update (if CS and PS are coordinated)
\end{enumerate}

\textbf{Types:}

\begin{itemize}
\tightlist
\item
  \textbf{Intra-SGSN RAU:} New RA served by same SGSN — local procedure, no GTP tunnel modification needed
\item
  \textbf{Inter-SGSN RAU:} New RA served by different SGSN — requires:

  \begin{enumerate}
  \def\labelenumi{\arabic{enumi}.}
  \tightlist
  \item
    New SGSN contacts old SGSN to get PDP contexts
  \item
    New SGSN informs GGSN to update tunnel endpoints
  \item
    New SGSN updates HLR with new location
  \item
    Old SGSN releases resources
  \end{enumerate}
\end{itemize}

\CNsubsection{9.5}{GPRS Attach Procedure}

Before any PDP context can be activated, the MS must \textbf{attach} to the GPRS network:

\begin{enumerate}
\def\labelenumi{\arabic{enumi}.}
\tightlist
\item
  MS sends \textbf{Attach Request} (with IMSI or P-TMSI)
\item
  SGSN authenticates the subscriber (same Ki-based mechanism as GSM)
\item
  SGSN updates HLR with subscriber’s new SGSN location
\item
  SGSN assigns a \textbf{P-TMSI} (Packet-TMSI — equivalent of TMSI for packet domain)
\item
  SGSN sends \textbf{Attach Accept} to MS
\end{enumerate}

After attach, the MS is ``registered'' for packet service but has no active data session yet — that requires PDP context activation.

\Needspace{14\baselineskip}

\CNsection{10}{EDGE — Enhanced Data Rates for GSM Evolution}

\CNchained\CNsubsection{10.1}{WHY EDGE Was Needed}

GPRS improved efficiency (packet switching) but was still limited by \textbf{GSM’s modulation scheme}:

\begin{itemize}
\tightlist
\item
  GSM/\allowbreak{}GPRS uses \textbf{GMSK} (Gaussian Minimum Shift Keying)
\item
  GMSK encodes \textbf{1~bit per symbol}
\item
  Maximum raw data rate per timeslot: \textasciitilde{}21.4~kbps (after coding)
\item
  Maximum practical throughput: \textasciitilde{}40-60~kbps with 4 timeslots
\end{itemize}

For emerging applications (email attachments, early mobile web), this wasn’t enough. But operators couldn’t yet afford to deploy entirely new 3G networks.

\textbf{EDGE’s value proposition:} Significantly more speed using the \textbf{same spectrum, same cell sites, same network} — just upgraded software and transceivers.

\Needspace{11\baselineskip}

\Needspace{22\baselineskip}

\CNsubsection{10.2}{WHAT EDGE Changes}

\CNchained\CNsubsubsection{Modulation: GMSK \ensuremath{\to} 8PSK}

{\def\LTcaptype{none} % do not increment counter
\Needspace{16\baselineskip}
\begin{longtable}[c]{@{}ccc@{}}
\toprule\noalign{}
\rowcolor{CNink}\CNth{Parameter} & \CNth{GPRS (GMSK)} & \CNth{EDGE (8PSK)} \\
\CNheadrulesep
\endhead
\bottomrule\noalign{}
\endlastfoot
Modulation & GMSK & 8PSK (8-Phase Shift Keying) \\
Bits per symbol & 1 & 3 \\
Raw bit rate per timeslot & 22.8~kbps & 69.2~kbps (3\ensuremath{\times} faster) \\
Symbol rate & 270.833 ksps & 270.833 ksps (unchanged!) \\
Carrier bandwidth & 200~kHz & 200~kHz (unchanged!) \\
Peak data rate (8 TS) & 171.2~kbps & 473.6~kbps \\
Practical peak & \textasciitilde{}80~kbps & \textasciitilde{}200-236~kbps \\
\end{longtable}
}

\begin{CNinsight}

\textbf{Key insight:} EDGE achieves 3\ensuremath{\times} the raw bit rate because 8PSK encodes 3~bits in every symbol instead of 1. The symbol rate and carrier bandwidth don’t change — it’s purely a modulation efficiency gain.

\end{CNinsight}

\CNsubsubsection{Trade-off: Why Not Always Use 8PSK?}

8PSK has \textbf{less distance between constellation points} than GMSK. This means:

\begin{itemize}
\tightlist
\item
  More susceptible to noise and interference
\item
  Higher bit error rate at cell edges
\item
  Not as robust in poor radio conditions
\end{itemize}

\textbf{EDGE’s solution: Link Adaptation}

\begin{itemize}
\tightlist
\item
  The system measures channel quality in real-time
\item
  Near the BTS (good signal): Use 8PSK with light coding (MCS-9: 59.2 kbps/\allowbreak{}slot)
\item
  Far from BTS (weak signal): Fall back to GMSK with heavy coding (MCS-1: 8.8 kbps/\allowbreak{}slot)
\item
  9 different Modulation and Coding Schemes (MCS-1 through MCS-9)
\end{itemize}

\Needspace{24\baselineskip}

\CNsubsection{10.3}{EDGE Coding Schemes}

{\def\LTcaptype{none} % do not increment counter
\Needspace{20\baselineskip}
\begin{longtable}[c]{@{}ccccc@{}}
\toprule\noalign{}
\rowcolor{CNink}\CNth{MCS} & \CNth{Modulation} & \CNth{Code Rate} & \CNth{Data Rate per Slot} & \CNth{When Used} \\
\CNheadrulesep
\endhead
\bottomrule\noalign{}
\endlastfoot
MCS-1 & GMSK & \textasciitilde{}0.53 & 8.8~kbps & Very poor channel \\
MCS-2 & GMSK & \textasciitilde{}0.66 & 11.2~kbps & Poor channel \\
MCS-3 & GMSK & \textasciitilde{}0.80 & 14.8~kbps & Below average \\
MCS-4 & GMSK & \textasciitilde{}1.0 & 17.6~kbps & Average (no redundancy) \\
MCS-5 & 8PSK & \textasciitilde{}0.37 & 22.4~kbps & Average channel \\
MCS-6 & 8PSK & \textasciitilde{}0.49 & 29.6~kbps & Good channel \\
MCS-7 & 8PSK & \textasciitilde{}0.76 & 44.8~kbps & Very good channel \\
MCS-8 & 8PSK & \textasciitilde{}0.92 & 54.4~kbps & Excellent channel \\
MCS-9 & 8PSK & \textasciitilde{}1.0 & 59.2~kbps & Best conditions \\
\end{longtable}
}

\CNsubsection{10.4}{EGPRS: EDGE Applied to Packet Data}

When EDGE modulation is applied to GPRS, the result is called \textbf{EGPRS (Enhanced GPRS)}:

\begin{itemize}
\tightlist
\item
  Same GPRS architecture (SGSN, GGSN, PDP contexts)
\item
  Same protocol stack
\item
  Only the physical layer and RLC layer are enhanced
\item
  Adds \textbf{Incremental Redundancy (IR):} If a block fails, retransmit with different puncturing; receiver combines attempts for better decoding probability
\end{itemize}

\Needspace{17\baselineskip}

\CNsubsection{10.5}{Network Changes Required for EDGE}

{\def\LTcaptype{none} % do not increment counter
\Needspace{13\baselineskip}
\begin{longtable}[c]{@{}
  >{\raggedright\arraybackslash\hspace{0pt}}m{(\linewidth - 4\tabcolsep) * \real{0.2963}}
  >{\raggedright\arraybackslash\hspace{0pt}}m{(\linewidth - 4\tabcolsep) * \real{0.4815}}
  >{\raggedright\arraybackslash\hspace{0pt}}m{(\linewidth - 4\tabcolsep) * \real{0.2222}}@{}}
\toprule\noalign{}
\rowcolor{CNink}\CNth{Change} & \CNth{Description} & \CNth{Cost} \\
\CNheadrulesep
\endhead
\bottomrule\noalign{}
\endlastfoot
New transceivers & BTS needs 8PSK-capable transmitters (linear amplifiers) & Moderate \\
Software upgrade & BSC, SGSN need new coding scheme support & Low \\
No new spectrum & Uses existing GSM frequencies & Zero \\
No new sites & Uses existing cell sites & Zero \\
No new core nodes & Same SGSN/\allowbreak{}GGSN infrastructure & Zero \\
\end{longtable}
}

This made EDGE extremely attractive as a ``2.5G+'' stepping stone.

\Needspace{14\baselineskip}

\CNkeepfig{m01-evolution}{8}

\CNsection{11}{Evolution Path: GSM \ensuremath{\to} GPRS \ensuremath{\to} EDGE \ensuremath{\to} 3G}

\CNchained\CNsubsection{11.1}{The Evolution Timeline}

\CNfigure{m01-evolution}{Evolution timeline from GSM to UMTS}

\Needspace{15\baselineskip}

\CNsubsection{11.2}{What Each Step Added}

{\def\LTcaptype{none} % do not increment counter
\Needspace{11\baselineskip}
\begin{longtable}[c]{@{}
  >{\raggedright\arraybackslash\hspace{0pt}}m{(\linewidth - 6\tabcolsep) * \real{0.1746}}
  >{\raggedright\arraybackslash\hspace{0pt}}m{(\linewidth - 6\tabcolsep) * \real{0.2381}}
  >{\raggedright\arraybackslash\hspace{0pt}}m{(\linewidth - 6\tabcolsep) * \real{0.2540}}
  >{\raggedright\arraybackslash\hspace{0pt}}m{(\linewidth - 6\tabcolsep) * \real{0.3333}}@{}}
\toprule\noalign{}
\rowcolor{CNink}\CNth{Generation} & \CNth{Key Innovation} & \CNth{Problem Solved} & \CNth{What Stayed the Same} \\
\CNheadrulesep
\endhead
\bottomrule\noalign{}
\endlastfoot
\textbf{GSM} & Digital cellular, roaming & Incompatible analog systems & — \\
\textbf{GPRS} & Packet switching overlay & Data was too expensive (circuit-switched) & GSM radio, BSS, spectrum \\
\textbf{EDGE} & 8PSK modulation + link adaptation & GPRS was too slow for emerging apps & GPRS architecture, spectrum, sites \\
\textbf{UMTS (3G)} & New air interface (WCDMA), 5~MHz carriers & EDGE still couldn’t support video/\allowbreak{}rich media & Core network concepts (SGSN \ensuremath{\to} modernized) \\
\end{longtable}
}

\CNsubsection{11.3}{WHY This Evolution Path Matters}

The GSM \ensuremath{\to} GPRS \ensuremath{\to} EDGE progression demonstrates a key principle in telecom engineering:

\begin{CNquote}

\textbf{Maximize return from existing infrastructure before deploying new generations.}

\end{CNquote}

Each step was designed to:

\begin{enumerate}
\def\labelenumi{\arabic{enumi}.}
\tightlist
\item
  \textbf{Reuse} existing investments (spectrum, sites, core network)
\item
  \textbf{Add} capability incrementally
\item
  \textbf{Allow} operators to upgrade gradually (not ``big bang'')
\end{enumerate}

This philosophy continues today: 4G \ensuremath{\to} 4G-Advanced \ensuremath{\to} 5G NSA \ensuremath{\to} 5G SA.

\Needspace{24\baselineskip}

\CNsection{}{Summary Table}

{\def\LTcaptype{none} % do not increment counter
\Needspace{18\baselineskip}
\begin{longtable}[c]{@{}
  >{\raggedright\arraybackslash\hspace{0pt}}m{(\linewidth - 6\tabcolsep) * \real{0.2917}}
  >{\raggedright\arraybackslash\hspace{0pt}}m{(\linewidth - 6\tabcolsep) * \real{0.2083}}
  >{\raggedright\arraybackslash\hspace{0pt}}m{(\linewidth - 6\tabcolsep) * \real{0.2500}}
  >{\raggedright\arraybackslash\hspace{0pt}}m{(\linewidth - 6\tabcolsep) * \real{0.2500}}@{}}
\toprule\noalign{}
\rowcolor{CNink}\CNth{Topic} & \CNth{GSM} & \CNth{GPRS} & \CNth{EDGE} \\
\CNheadrulesep
\endhead
\bottomrule\noalign{}
\endlastfoot
\textbf{Switching} & Circuit & Packet & Packet \\
\textbf{Modulation} & GMSK & GMSK & GMSK + 8PSK \\
\textbf{Peak Data Rate} & 9.6~kbps (14.4 CSD) & 171.2~kbps (theoretical) & 473.6~kbps (theoretical) \\
\textbf{Practical Data} & 9.6~kbps & 30–80~kbps & 120–236~kbps \\
\textbf{Billing} & Per minute & Per MB & Per MB \\
\textbf{Always-on?} & No (dial-up) & Yes & Yes \\
\textbf{New Infrastructure} & Complete network & + SGSN, GGSN, PCU & + 8PSK transceivers \\
\textbf{New Spectrum?} & Yes (900/\allowbreak{}1800~MHz) & No & No \\
\end{longtable}
}

\CNboxsection{Review Questions}

\begin{CNbox}[unbreakable]{q}{Review Questions}

\begin{enumerate}
\def\labelenumi{\arabic{enumi}.}
\tightlist
\item
  Explain why the BSS is designed with ``dumb'' BTSs and an intelligent BSC. What are the advantages of this split?
\item
  Why does the VLR exist? What would happen if every operation required a query to the HLR?
\item
  In the GSM authentication process, why does the Ki never leave the SIM or AuC? What would be the consequence if it did?
\item
  Explain why circuit switching is fundamentally inefficient for bursty data traffic like web browsing.
\item
  What is the role of GTP tunneling between SGSN and GGSN? Why can’t packets just be routed normally?
\item
  Why does EDGE fall back to GMSK in poor radio conditions instead of always using 8PSK?
\item
  If EDGE can achieve \textasciitilde{}236~kbps, why was 3G (UMTS) still needed?
\end{enumerate}

\end{CNbox}

\CNsection{}{References}

\begin{itemize}
\tightlist
\item
  3GPP TS 03.02: Network Architecture (GSM)
\item
  3GPP TS 23.060: GPRS Service Description (Stage 2)
\item
  3GPP TS 43.064: GERAN Overall Description (EDGE)
\item
  3GPP TS 45.001: Physical Layer on the Radio Path
\item
  Halonen, Romero, Melero: ``GSM, GPRS and EDGE Performance'' (Wiley, 2003)
\item
  Eberspächer, Vögel, Bettstetter: ``GSM – Architecture, Protocols and Services'' (Wiley, 2009)
\end{itemize}
